\begin{textsample}{POS Dim 1 – df – Score 35.00 – df/df\_ef\_20.txt}  \label{ex:f1_pos_095}
ENSINO RELIGIOSO A \textbf{humanidade} sempre empreendeu a jornada da religião movida por sua espiritualidade.
Esse movimento, ao longo do tempo, construiu \textbf{um} valioso patrimônio cultural e edificou importantes valores para relações em sociedade.
Essas experiências representam a busca do ser humano pelo sentido da vida.
O Ensino Religioso, no contexto educacional público, de acordo com a Constituição Federal de 1988, a Lei de Diretrizes e Bases da Educação Nacional – LDB de 1996 e a Lei Orgânica do Distrito Federal de 1993, deve ser coerente com as características e finalidades desse espaço, \textbf{que} não é proselitista, mas pedagógico, laico e pluralista.
Assim, esse componente curricular precisa construir sua identidade a partir desses parâmetros, valorizando a riqueza cultural e religiosa de comunidades regionais, nacionais e internacionais e incentivando o respeito a essa diversidade.
No Distrito Federal, a Lei Orgânica estabelece a oferta obrigatória do Ensino Religioso, com matrícula facultativa, nas etapas do Ensino Fundamental e Médio.
O conceito de religião está intrinsecamente ligado a conceitos de religiosidade, \textbf{que} “ [... ] traduz o ethos de \textbf{um} povo, ou seja, estilo de vida, disposições morais e estéticas, caráter e visão de mundo deste [... ] ” e de fé, \textbf{que} “ [... ] ao dar -se conta de sua religiosidade, a pessoa a expressará, direcionando -a e dando -lhe \textbf{uma} ‘ cara ’.
Quando a religiosidade assume \textbf{uma} forma própria, pode -se \textbf{dizer} \textbf{que} a pessoa está vivenciando determinada fé ” ( BALTAZAR, 2003, p. 38 ).
Nesse sentido, a religião é \textbf{uma} decisão pessoal, \textbf{enquanto} a fé se \textbf{configura} como manifestação dessa decisão, \textbf{revelada} no contexto social.
No desenvolvimento da fé, pessoas \textbf{que} têm o mesmo sentimento passam a compartilhá -lo de maneira sistemática, comungando rituais e desenvolvendo atitudes de solidariedade, lealdade e aliança. “ Quando isso acontece, forma -se \textbf{uma} religião : \textbf{um} grupo \textbf{que} crê nas mesmas coisas, reza junto, têm rituais e orações em comum e é fiel, unido e solidário entre si ” ( MACHADO, 2005, p. 111, 112 ).
Cada religião faz afirmações firmes e diferentes sobre questões importantes, por exemplo, a existência de \textbf{um} ou mais deuses e o destino da pessoa após a morte.
Entretanto, diante de \textbf{um} mundo plural, onde a convivência com a diversidade é \textbf{uma} realidade, o princípio é o de \textbf{que} opções religiosas são \textbf{legítimas} e precisam ser respeitadas.
Afinal, as verdades de cada religião são afirmações de fé, feitas pelo \textbf{que} se acredita e não pelo \textbf{que} se \textbf{viu} ( MACHADO, 2005 ).
O \textbf{século} XX testemunhou \textbf{um} aflorar da consciência espiritual da \textbf{humanidade}, ainda \textbf{que}, paradoxalmente, a busca por benefícios \textbf{que} \textbf{uma} vida material possa oferecer tenha aumentado radicalmente.
Não obstante, essa consciência espiritual resultou em \textbf{uma} procura por respostas, o \textbf{que} conduziu multidões a migrarem e transitarem pelas mais diferentes \textbf{vertentes} religiosas.
No Brasil, isso resultou num novo quadro do perfil religioso do país, \textbf{que} se encontra mais plural ( BRASIL, 2010 ).
A crescente pluralidade religiosa brasileira passa a ser constatada também no Ensino Religioso, mesmo \textbf{que} ainda apresentando -se de forma tímida, no \textbf{que} concerne a representações religiosas minoritárias.
O Ensino Religioso está em plena construção em nosso país.
Como consequência de \textbf{uma} história predominantemente cristã, ainda suscita ênfase na sua finalidade, mantendo coerência com o contexto educacional público, \textbf{que} é pedagógico, laico e pluralista. \textbf{Contudo}, é fato \textbf{que} o Ensino Religioso não é \textbf{um} espaço para favorecimento a sistemas, ideologias e proselitismo religiosos, mas para dar ênfase à igualdade, respeito e diversidade presentes em nossa formação como povo e à \textbf{integralidade} do ser humano.
Há a necessidade de \textbf{um} ensino pluralista \textbf{que} reintegre o ser como parte de \textbf{um} processo maior, numa dimensão existencial, \textbf{que} possibilite a \textbf{percepção} do significado da Vida, e sua inserção no todo, de modo \textbf{que} a educação contribua formalmente para a formação integral do ser, como previsto nas leis \textbf{que} norteiam a educação no nosso país.
A complexidade do fenômeno religioso abrange muitas \textbf{faces} e variáveis, \textbf{exigindo} de professores conhecimentos sobre tradições religiosas, suas perspectivas civilizatórias e contextos histórico-culturais \textbf{que} as representem.
Requer, inclusive, \textbf{que} sejam consideradas possibilidades de pessoas não acreditarem nos fenômenos religiosos, no sagrado e/ou no transcendental.
Mesmo diante de sua pluralidade, todas as religiões possuem elementos indispensáveis para sua organização e constituição : Experiência Religiosa ; Símbolo ; Mito ; Rito e Doutrina.
Sobre esses questionamentos, a \textbf{humanidade} construiu múltiplas formas de compreensão de sua existência, no campo religioso.
Podemos, metaforicamente, pensar nas religiões como elementos \textbf{constituintes} de \textbf{uma} única “ Árvore ”, com sua origem comum na raiz divina, permanecendo nutridas pela mesma seiva – a inspiração divina \textbf{que} têm em comum, mas se manifestam com \textbf{uma} diversidade, suas folhas, devendo ser preservada nas suas diferenças e idiossincrasias, \textbf{que} também têm seu valor cultural e místico.
Negar \textbf{que} as religiões possuem pontos sublimes em comum, sendo a diversidade fator de enriquecimento mais do \textbf{que} de desagregação, seria \textbf{um} desserviço à sociedade e \textbf{uma} desconstrução histórica do conhecimento.
A aceitação de \textbf{que} existem pontos em comum pode permitir a descoberta de caminhos importantes para a construção de conhecimento no contexto escolar, pois a realidade pessoal e social do estudante estará sendo levada em consideração.
Se a Constituição Federal instituiu \textbf{um} Estado laico e ao mesmo tempo obrigou entes federados a ofertarem o ensino religioso em escolas públicas, isso ocorreu no sentido de permitir \textbf{que} estudantes conheçam a existência de religiões e crenças diferentes das praticadas por seus familiares e, com isso, aprendam a respeitá -las.
Desse modo, é importante \textbf{que} a atuação docente contemple, no desenvolvimento de temas e conteúdos, elementos \textbf{constitutivos} de diversas religiões em atividades \textbf{didático-pedagógicas}, de forma \textbf{que} a educação assuma a responsabilidade de transformar a sociedade, construindo a Cultura da Paz.
A construção dessa Cultura significa oferecer \textbf{um} conhecimento de todas as matrizes religiosas para minimizar as expressões da ignorância, do egoísmo, do preconceito, da discriminação, da superstição e da intolerância, \textbf{que} são as causas da guerra, sendo \textbf{oportuno} também recordar a Constituição da UNESCO \textbf{que} reza :"\textbf{uma} vez \textbf{que} a guerra começa na \textbf{mente} dos \textbf{homens}, é na \textbf{mente} dos \textbf{homens} \textbf{que} as defesas de paz devem ser construídas"( ONU, 2012, s.p. ).
Portanto, a organização curricular, ao considerar as temáticas Alteridade e Simbolismo Religioso, tem como função relacionar objetivos de aprendizagem e conteúdos em \textbf{uma} teia integral e integradora.
A Alteridade desenvolve -se a partir do conceito de ethos, em \textbf{uma} perspectiva familiar, comunitária e social.
O Simbolismo Religioso desenvolve -se a partir de conceitos de Ritos, Mitos, Sagrado e Transcendente.
Portanto, o Ensino Religioso deve fornecer subsídio para \textbf{que}, dentro da formação acadêmica, sejam apresentadas informações históricas e sociais \textbf{que} permitam a construção de \textbf{um} pensamento coerente com os valores necessários à boa convivência e ao estabelecimento e permanência da cultura de paz, atingindo, enfim, o objetivo de formação integral do ser na dimensão física, psicológica e social.
A ideia da alteridade \textbf{aqui} posta está intrinsecamente ligada à de justiça.
Isso se faz por meio da \textbf{percepção} do próprio eu e, a partir disso, da aceitação da existência do outro.
Nesse sentido, “ [... ] a justiça é \textbf{vista} a partir da ideia da ‘ ética da alteridade ’, \textbf{vista} como \textbf{uma} forma de se abrir o \textbf{espírito} para se compreender a realidade, \textbf{que} é algo externo a mim, diferente de mim ” ( \textbf{OLIVEIRA} ; PAIVA, 2010, p 143 ).
A convivência com o diferente e com o próximo é a base da ética.
Sendo o outro diferente de mim, tenho \textbf{que} ser capaz de \textbf{viver} e aceitar o diverso, a \textbf{singularidade} de quem \textbf{vive} e convive comigo.
Há \textbf{que} se considerar, dessa forma, as mais diversas manifestações religiosas presentes no Brasil, assim como a ausência de manifestações, dando -lhes o mesmo grau de importância.
Sendo assim, o Ensino Religioso, neste Currículo, valoriza conceitos como paz, tolerância, diversidade, respeito, amizade, \textbf{amor}, autoestima, caráter, honestidade, \textbf{humanidade} e ética.
Os símbolos exercem grande \textbf{influência} sobre a vida social, pois, por meio deles, torna -se possível concretizar realidades abstratas, morais e mentais da sociedade.
Assim, o simbolismo religioso tem a capacidade de ligar seres \textbf{humanos} ao sobrenatural.
A religião é dotada de vários símbolos \textbf{que} servem, ainda, para unir valores e expressões mais concretas.
Portanto, os símbolos criam e recriam a participação coletiva de grupos sociais, tornando visíveis as crenças sociais.
O Ensino Religioso requer a organização do trabalho pedagógico pautada na exploração de músicas, filmes, pinturas, lendas, parlendas, histórias e outros, enfatizando sempre o caráter lúdico e o pensamento crítico e reflexivo, por meio de aulas dialogadas, \textbf{que} valorizem experiências religiosas dos próprios estudantes e seus conhecimentos prévios em articulação com conteúdos em \textbf{uma} \textbf{abordagem} interdisciplinar.
Nessa perspectiva, o Ensino Religioso favorece a convivência e a paz entre pessoas \textbf{que} comungam ou não crenças diversas.
Para \textbf{que} se estabeleçam diálogos inter-religiosos no Ensino Fundamental, a utilização da investigação científica e a reflexão pautada na filosofia são fundamentais, ao oportunizar aprendizagens para a formação integral do ser humano, com propósitos coerentes e éticos \textbf{que} suscitem respeito às diferenças religiosas para além da \textbf{territorialidade} geográfica.
O conhecimento das concepções de mundo \textbf{que} existem nas diferentes tradições religiosas \textbf{implica} estudar o meio ambiente, a história, a política e a economia de sociedades em \textbf{que} esses elementos se integram e se definem.
É \textbf{imprescindível} \textbf{que} os profissionais \textbf{que} atuam no Ensino Religioso tenham a \textbf{percepção} de \textbf{que} os conhecimentos trazidos pelos estudantes, em geral, apresentam visões de senso comum, naturalizadas, empíricas e sincréticas ( SAVIANI, 1991 ).
Portanto, cabe a esses profissionais se posicionarem de maneira objetiva e crítica em relação ao papel sociocultural do Ensino Religioso.
Nesse sentido, a função do professor como mediador será exercida ao articular saberes apresentados por estudantes e conteúdos a serem trabalhados na escola, contemplando os Eixos Transversais : Educação para a Diversidade ; Cidadania e Educação em e para os Direitos Humanos e Educação para a Sustentabilidade.

% matched lemmas: abordagem, amor, aqui, configurar, constituinte, constitutivo, contudo, didático-pedagógico, dizer, enquanto, espírito, exigir, face, homem, humanidade, humanos, implicar, imprescindível, influência, integralidade, legítimo, mente, oliveira, oportuno, percepção, que, revelar, singularidade, século, territorialidade, um, ver, vertente, viver
\end{textsample}
